\BourbakiExercisesSection

\begin{enumerate}
\def\labelenumi{\arabic{enumi}.}
\tightlist
\item
  确定一个具有 \(n\) 个元素的集合上的所有群结构，对 \(2 \leqslant n \leqslant 6\) （参见 § 2，习题5）。确定这些群的子群与商群，以及它们的若尔当–赫尔德序列。
\end{enumerate}

¶ 2. (a) 集合 E 上的一个结合律 \((x,y)\mapsto xy\) 是群律，如果存在 \(e\in\mathrm{E}\) 使得对所有 \(x\in\mathrm{E}\) 有 \(ex=x\) ，且如果对每个 \(x\in\mathrm{E}\) 存在 \(x'\) 使得 \(x^{\prime}x=e\) （证明 \(xx^{\prime}=e\) ，考察复合 \(x^{\prime}xx^{\prime}\) ；由此推出 \(e\) 是单位元）。

\begin{enumerate}
\def\labelenumi{(\alph{enumi})}
\setcounter{enumi}{1}
\tightlist
\item
  证明如果对所有 \(x \in E\) ，左平移 \(\gamma_{x}\) 是 E 到 E 上的一个映射，并且如果存在某个 \(a \in E\) 使得右平移 \(\delta_{a}\) 是 E 到 E 上的一个映射，则同样的结论成立。
\end{enumerate}

\begin{enumerate}
\def\labelenumi{\arabic{enumi}.}
\setcounter{enumi}{2}
\item
  在群 G 中，每个有限非空稳定子集 H 都是 G 的子群（参见 § 2，习题 6）。
\item
  设 A 和 B 是群 G 的两个子群。
\end{enumerate}

\begin{enumerate}
\def\labelenumi{(\alph{enumi})}
\item
  证明包含 A 和 B 的最小子群（即由 \(\mathbf{A} \cup \mathbf{B}\) 生成的子群）与奇数个（任取的）元素的序列 \((x_{i})_{1 \leqslant i \leqslant 2n + 1}\) 的复合所成的集合相同，其中 \(x_{i} \in \mathbf{A}\) 对 \(i\) 为奇数， \(x_{i} \in \mathbf{B}\) 对 \(i\) 为偶数。
\item
  为使AB成为G的子群（此时它是由A∪B生成的子群），其充分必要条件为A与B可交换，即AB = BA。
\item
  若 A 与 B 可交换，且 C 是包含 A 的子群，则 A 与 \(B \cap C\) 可交换，且 \(\mathbf{A}(\mathbf{B} \cap \mathbf{C}) = \mathbf{C} \cap (\mathbf{AB})\) 。
\end{enumerate}

\begin{enumerate}
\def\labelenumi{\arabic{enumi}.}
\setcounter{enumi}{4}
\item
  若群G的一个子群的指数为2，则它在G中正规。
\item
  设 \((\mathbf{G}_{\alpha})\) 是群 \(\mathbf{G}\) 的一族正规子群，使得 \(\bigcap_{\alpha} \mathbf{G}_{\alpha} = \{e\}\) ；证明 \(\mathbf{G}\) 同构于乘积群 \(\prod_{\alpha} (\mathbf{G} / \mathbf{G}_{\alpha})\) 的一个子群。
\item
  若 G 是两个子群 A 与 B 的直积，且 H 是 G 的一个子群，使得 \(A \subset H\) ，则 H 是 A 与 \(H \cap B\) 的直积。
\item
  设 A 与 A' 是两个群，且 G 是 A × A' 的一个子群。我们记：
\end{enumerate}

\[
\mathrm{N} = \mathrm{G} \cap (\mathrm{A} \times \{e \}), \quad \mathrm{H} = \operatorname{pr} _ {1} (\mathrm{G})
\]

\[
\mathrm{N} ^ {\prime} = \mathrm{G} \cap (\{e \} \times \mathrm{A} ^ {\prime}), \quad \mathrm{H} ^ {\prime} = \operatorname{pr} _ {2} (\mathrm{G}).
\]

\begin{enumerate}
\def\labelenumi{(\alph{enumi})}
\tightlist
\item
  证明 N 在 H 中正规，且 N' 在 H' 中正规；定义同构
\end{enumerate}

\[
\mathrm{H} / \mathrm{N} \rightarrow \mathrm{G} / (\mathrm{N} \times \mathrm{N} ^ {\prime}) \rightarrow \mathrm{H} ^ {\prime} / \mathrm{N} ^ {\prime}.
\]

\begin{enumerate}
\def\labelenumi{(\alph{enumi})}
\setcounter{enumi}{1}
\tightlist
\item
  假设 H = A， \(H' = A'\) ，群 A 与 \(A'\) 具有有限长度，且 A 的若尔当–赫尔德序列的任一商群都不同构于 \(A'\) 的若尔当–赫尔德序列的任一商群。证明 \(G = A \times A'\) 。
\end{enumerate}

\begin{enumerate}
\def\labelenumi{\arabic{enumi}.}
\setcounter{enumi}{8}
\item
  设 G 是一个不只由单位元组成的群。假设存在 G 的一个有限子集 S，使得 G 由 S（相应地，由元素 \(gsg^{-1}, g \in G, s \in S\) ）生成。设 N 为 G 的异于 G 的子群（相应地，正规子群）所成的集合。证明 N 在包含序下是归纳的。由此推出存在 G 的一个正规子群 H，使得 G/H 是单群。
\item
  设 H 是群 G 的一个正规子群，且包含在 G 的中心内。证明：若 G/H 是单生成群，则 G 是交换群。
\item
  如果群 G 中除单位元外的所有元素阶均为 2，则 G 是交换的；若 G 有限，则其阶 n 是 2 的幂（通过对 n 进行归纳论证）。
\item
  设 G 是一个群，使得对一个固定的整数 \(n > 1\) ， \((xy)^n = x^n y^n\) 对所有 \(x \in G\) 、 \(y \in G\) 成立。若 \(G^{(n)}\) 表示由 \(x^n\) （其中 \(x\) 遍历 G）组成的集合，而 \(G_{(n)}\) 表示由 \(x \in G\) （满足 \(x^n = e\) ）组成的集合，证明 \(G^{(n)}\) 与 \(G_{(n)}\) 是 G 的正规子群；若 G 有限， \(G^{(n)}\) 的阶等于 \(G_{(n)}\) 的指数。证明对所有 G 中的 \(x, y\) ，也有 \(x^{1-n}y^{1-n} = (xy)^{1-n}\) ，并由此推出 \(x^{n-1}y^n = y^n x^{n-1}\) ；由此得出结论：G 中形如 \(x^{n(n-1)}\) 的元素所成的集合生成 G 的一个交换子群。
\item
  设 G 是一个群。若 S 是一个单群，则称 S 出现在 G 中，如果存在 G 的两个子群 H、H'，其中 H 是 H' 的正规子群，使得 H'/H 同构于 S。令 in(G) 表示出现在 G 中的单群的同构类集合。
\end{enumerate}

\begin{enumerate}
\def\labelenumi{(\alph{enumi})}
\item
  证明 in(G) = ∅ ⇔ G = \{e\}.
\item
  若 \(\mathbf{H}\) 是 \(\mathbf{G}\) 的子群，证明 \(\operatorname{in}(\mathbf{H}) \subset \operatorname{in}(\mathbf{G})\) ；此外，若 \(\mathbf{H}\) 是正规的，则
\end{enumerate}

\[
\operatorname{in} (\mathrm{G}) = \operatorname{in} (\mathrm{H}) \cup \operatorname{in} (\mathrm{G} / \mathrm{H}).
\]

\begin{enumerate}
\def\labelenumi{(\alph{enumi})}
\setcounter{enumi}{2}
\item
  设 \(\mathbf{G}_1\) 和 \(\mathbf{G}_2\) 是两个群。证明以下两个性质的等价性：
\item
  \(\operatorname{in}(\mathbf{G}_1) \cap \operatorname{in}(\mathbf{G}_2) = \varnothing\) ;
\end{enumerate}

\begin{enumerate}
\def\labelenumi{(\roman{enumi})}
\setcounter{enumi}{1}
\tightlist
\item
  \(\mathbf{G}_1\times \mathbf{G}_2\) 的每个子群都具有形式 \(\mathbf{H}_1\times \mathbf{H}_2\) ，其中 \(\mathbf{H}_1\subset \mathbf{G}_1\) ， \(\mathbf{H}_2\subset \mathbf{G}_2\) 。
\end{enumerate}

（使用习题8。）

*当 \(\mathbf{G}_1\) 和 \(\mathbf{G}_2\) 是有限群时，证明（i）和（ii）等价于：

\begin{enumerate}
\def\labelenumi{(\roman{enumi})}
\setcounter{enumi}{2}
\tightlist
\item
  \(\operatorname{Card}(\mathbf{G}_1)\) 和 \(\operatorname{Card}(\mathbf{G}_2)\) 是互素的。*
\end{enumerate}

\begin{enumerate}
\def\labelenumi{\arabic{enumi}.}
\setcounter{enumi}{13}
\tightlist
\item
  设G为一个群，H为G的一个子群。G的任意子集T，若它与模H的每个左陪集恰好交于一点，则称T为模H的左陪集的一个代表系（参见集合论，II，§6，第2号）；此条件等价于以下条件：
\end{enumerate}

映射 \((x,y)\mapsto xy\) 是 \(\mathbf{T}\times \mathbf{H}\) 到 \(\mathbf{G}\) 上的一个双射。

\begin{enumerate}
\def\labelenumi{(\alph{enumi})}
\item
  设 \(\mathbf{T}\) 是这样的一个代表系；对所有 \(g \in \mathbf{G}\) 与所有 \(t \in \mathbf{T}\) ，设 \(x(g, t) \in \mathbf{T}\) 与 \(y(g, t) \in \mathbf{H}\) 满足 \(gt = x(g, t)y(g, t)\) 。设 \(\mathbf{S}\) 是 \(\mathbf{G}\) 的一个生成 \(\mathbf{G}\) 的子集。证明元素 \(y(g, t), g \in \mathbf{S}, t \in \mathbf{T}\) 生成 \(\mathbf{H}\) 。（设 \(\mathbf{H}'\) 为 \(\mathbf{H}\) 的由这些元素生成的子群，证明 \(\mathbf{T}.\mathbf{H}'\) 在诸平移 \(\boldsymbol{\gamma}_g, g \in \mathbf{S}\) 下稳定，从而在诸 \(\boldsymbol{\gamma}_g, g \in \mathbf{G}\) 下也稳定，并由此推出 \(\mathbf{T}.\mathbf{H}' = \mathbf{G}\) ，从而 \(\mathbf{H}' = \mathbf{H}\) 。）
\item
  假设 (G:H) 是有限的。证明：G 能被有限子集生成当且仅当 H 能被有限子集生成。
\end{enumerate}

¶ 15. 设F是带算子群G的一组稳定子群；若F的每个子集在包含序下都有极大（相应地，极小）元，则称F满足极大条件（相应地，极小条件）。

假设群G的所有稳定子群构成的集合满足极小条件。

\begin{enumerate}
\def\labelenumi{(\alph{enumi})}
\item
  证明不存在 G 的同构于 G 且异于 G 的稳定子群（用反证法论证：表明该假设将蕴含存在 G 的稳定子群的一个严格递减无限序列）。
\item
  G 的不只由 e 组成的正规稳定子群所成集合的极小元素称为极小正规稳定子群。设 M 是 G 的一族极小正规稳定子群，S 是 G 的包含 M 中所有子群的最小稳定子群；证明 S 是 G 的有限个极小正规稳定子群的直积（令 \((\mathbf{M}_{n})\) 为 G 的属于 M 且满足下述条件的一列极小正规稳定子群： \(M_{n+1}\) 不包含在由 \(M_{1}, M_{2}, \ldots, M_{n}\) 的并集生成的稳定子群中；令 \(S_{k}\) 为由诸 \(M_{n}\) （指标 \(n \geqslant k\) ）的并集生成的稳定子群；证明从某个秩开始 \(S_{k+1} = S_{k}\) ，因而 \((\mathbf{M}_{n})\) 是一个有限序列；最后使用第 9 号，命题 15）。
\item
  如果G是一个无算子的群，证明G的每个极小正规子群M是有限个彼此同构的单群的直积（设N是M的极小正规子群；证明M是G中包含所有子群 \(aNa^{-1}\) （其中a遍历G）的最小子群，然后将(b)应用于群M）。
\end{enumerate}

\begin{enumerate}
\def\labelenumi{\arabic{enumi}.}
\setcounter{enumi}{15}
\tightlist
\item
  如果带算子的群 G 的稳定子群集合满足极大条件和极小条件（习题 15），则 G 有一个若尔当–赫尔德序列（对 G 的一个子群 H，考虑 H 的异于 H 的正规稳定子群所成集合的一个极大元素）。
\end{enumerate}

¶ 17. 设G是一个带算子的群；G的一个合成列（G \(_{i}\) ）称为正规的，如果所有G \(_{i}\) 都是G的正规稳定子群；一个正规列Σ称为主列，如果它是严格递减的，并且不存在异于Σ、比Σ更细且严格递减的正规列。

\begin{enumerate}
\def\labelenumi{(\alph{enumi})}
\item
  若 \((\mathbf{G}_{i})\) 和 \((\mathbf{H}_{j})\) 是 G 的两个正规列，证明存在两个等价的正规列，分别比 \((\mathbf{G}_{i})\) 和 \((\mathbf{H}_{j})\) 更细（通过考虑 G 的一个合适的算子域来应用 Schreier 定理）。通过将子群 \(G_{ij}^{\prime} = G_{i} \cap (G_{i+1}H_{j})\) 和 \(H_{ji}^{\prime} = H_{j} \cap (H_{j+1}G_{i})\) “插入”到序列 \((\mathbf{G}_{i})\) 和 \((\mathbf{H}_{j})\) 中，给出这个命题的第二个证明。
\item
  如果G有主列，则G的任意两个主列等价；对于每个严格递减的正规列 \(\Sigma\) ，存在一个比 \(\Sigma\) 更细的主列。由此推出，G有主列的充分必要条件是G的正规稳定子群集合满足极大条件和极小条件。
\item
  如果G是一个无算子的群，有主列，并且G的子群集合满足极小条件，则每个商群
\end{enumerate}

\(G_{i}/G_{i+1}\) 是有限个彼此同构的单子群的直积（使用习题15）。

¶ 18. (a) 设 G 是一个带算子的群，由一族 G 的单正规稳定子群 \((\mathrm{H}_{i})_{i\in\mathbf{I}}\) 的并集生成。证明存在 I 的一个子集 J，使得 G 是族 \((\mathrm{H}_{i})_{i\in\mathbf{J}}\) 的限制和（将 Zorn 引理应用于 I 的具有下述性质的子集 K 所成的集合：由 \(H_{i}\) （ \(i\in K\) ）的并集生成的子群是该族的限制和）。

\begin{enumerate}
\def\labelenumi{(\alph{enumi})}
\setcounter{enumi}{1}
\tightlist
\item
  设 A 是 G 的一个正规稳定子群。证明存在 I 的一个子集 \(J_{A}\) ，使得 G 是 A 与诸 \(H_{i}\) （ \(i \in J_{A}\) ）的限制和。由此推出 A 同构于诸 \(H_{i}\) 的一个子族的限制和。
\end{enumerate}

¶ 19. (a) 设 G 是一个群，使得 G 的每个异于 G 的正规子群都包含在 G 的一个异于 G 的直因子中。证明 G 是一族单子群的限制和。（设 K 为 G 的由所有单子群的并集生成的子群。假设 K ≠ G；若 x \ensuremath{\in} G−K，考虑 G 的一个正规子群 M，使其在包含 K 而不包含 x 的那些正规子群中是极大的。若 G = M' × S，其中 M' 包含 M 且 S ≠ \{e\} ，证明 x ∉ M' ，从而 M' = M；另一方面，若 N 是 S 的一个正规子群，证明 M × N = G 并由此得出 S 是单群，从而得到矛盾。借助习题 18 得出结论。）得出逆命题（参见习题 18）。

\begin{enumerate}
\def\labelenumi{(\alph{enumi})}
\setcounter{enumi}{1}
\tightlist
\item
  为使在群 G 中存在一族 \((\mathrm{N}_{\alpha})\) 正规单子群，使得诸 \(G/N_{\alpha}\) 都是单群且 \(\bigcap_{\alpha} N_{\alpha} = \{e\}\) ，必要且充分的是：对 G 的每个正规子群 \(N \neq \{e\}\) ，存在一个正规子群 \(N' \neq G\) ，使得 \(NN' = G\) 。（要看出该条件是必要的，考虑一个异于 e 的 \(x \in N\) 以及一个不包含 x 的 \(N_{\alpha}\) 。要看出该条件是充分的，对 G 中所有 \(x \neq e\) ，考虑由 x 生成的正规子群 N 以及一个正规子群 \(N' \neq G\) ，使得 \(NN' = G\) 。证明：若 G 的正规子群 M 在包含 \(N'\) 而不包含 x 的 G 的正规子群中是极大的，则它在所有正规子群 \(\neq G\) 所成的集合中也是极大的，并由此得出结论。）
\end{enumerate}

\begin{enumerate}
\def\labelenumi{\arabic{enumi}.}
\setcounter{enumi}{19}
\tightlist
\item
  \begin{enumerate}
  \def\labelenumii{(\alph{enumii})}
  \tightlist
  \item
  设 G 是一个群，是两个子群 A 和 B 的直积。设 \(C_{A}\) 为 A 的中心，并设 N 是 G 的一个正规子群，使得 \(N \cap A = \{1\}\) 。证明 \(N \subset C_{A} \times B\) 。
  \end{enumerate}
\end{enumerate}

\begin{enumerate}
\def\labelenumi{(\alph{enumi})}
\setcounter{enumi}{1}
\tightlist
\item
  设 \((\mathbf{S}_i)_{i\in I}\) 是非交换单群的一个有限族。证明 \(\prod_{i\in I}S_i\) 的每个正规子群都等于诸部分积 \(\prod_{i\in J}S_i\) 之一，其中 \(J\) 是 \(I\) 的一个子集。
\end{enumerate}

\begin{enumerate}
\def\labelenumi{\arabic{enumi}.}
\setcounter{enumi}{20}
\tightlist
\item
  设 G 是一个有限长度的群。证明以下性质的等价性：
\end{enumerate}

\begin{enumerate}
\def\labelenumi{(\alph{enumi})}
\item
  G 是一些彼此同构的单群的乘积。
\item
  G 的任何异于 \(\{1\}\) 与 G 的子群在 G 的所有自同构下都不是稳定的。
\end{enumerate}

¶ 22. 设 E 是一个拟群（§ 3，习题 6），以乘法记号书写，具有一个幂等元 e，并且满足 \((xy)(zt) = (xz)(yt)\) 对所有 x, y, z, t 成立。设 \(u(x)\) 表示 E 中满足 \(u(x)e = x\) 的元素， \(v(x)\) 表示 E 中满足 \(ev(x) = x\) 的元素；证明合成律 \((x, y) \mapsto u(x)v(y)\) 是 E 上的一个交换群律，e 在其下是单位元，映射 \(x \mapsto xe\) 与 \(y \mapsto ey\) 是该群结构的可交换的自同态，并且 \(xy = u(xe)v(ey)\) （先建立恒等式 \(e(xy) = (ex)(ey)\) , \((xy)e = (xe)(ye)\) , \(e(xe) = (ex)e\) ；然后注意关系 x = y、ex = ey 与 xe = ye 是等价的）。得出逆命题。

¶ 23. 考虑集合 E 上的一个内律，它并非处处有定义，以乘法记号书写，并满足以下条件：

\begin{enumerate}
\def\labelenumi{(\arabic{enumi})}
\item
  如果合成 \((xy)z, x(yz)\) 之一有定义，则另一个也有定义，并且两者相等；
\item
  若 \(x, x', y\) 使得 \(xy\) 与 \(x'y\) （相应地 \(yx\) 与 \(yx'\) ）都有定义且相等，则 \(x = x'\) ；
\item
  对所有 \(x \in \mathbf{E}\) ，存在三个元素 \(e_x, e_x'\) 与 \(x^{-1}\) ，使得 \(e_xx = x\) , \(xe_x' = x\) , \(x^{-1}x = e_x'\) ； \(e_x\) 称为 \(x\) 的左单位元， \(e_x'\) 称为 \(x\) 的右单位元， \(x^{-1}\) 称为 \(x\) 的逆（这是语言上的一种滥用）。
\end{enumerate}

\begin{enumerate}
\def\labelenumi{(\alph{enumi})}
\item
  证明合成 \(xx^{-1}, x^{-1}e_x, e_x'x^{-1}, e_x e_x, e_x' e_x'\) 都有定义，并且 \(xx^{-1} = e_x, x^{-1}e_x = e_x'x^{-1} = x^{-1}, e_x e_x = e_x, e_x' e_x' = e_x'\) 。
\item
  每个幂等元 \(e\) （属于 \(\mathbf{E}\) ，§ 1，第 4 号）是所有 \(x\) （使 \(ex\) 有定义者）的左单位元，并且是所有 \(y\) （使 \(ye\) 有定义者）的右单位元。
\item
  为使合成 \(xy\) 有定义，必要且充分的是 \(x\) 的右单位元与 \(y\) 的左单位元相同（要看出该条件是充分的，利用关系式 \(e_y = yy^{-1}\) ）；若 \(xy = z\) ，则 \(x^{-1}z = y\) , \(zy^{-1} = x\) , \(y^{-1}x^{-1} = z^{-1}\) , \(z^{-1}x = y^{-1}\) , \(yz^{-1} = x^{-1}\) （这些关系式左边出现的合成都是有定义的）。
\item
  对任意两个幂等元 \(e, e'\) （属于 \(\mathbf{E}\) ），设 \(\mathbf{G}_{e, e'}\) 表示这样的 \(x \in \mathbf{E}\) 所成的集合： \(e\) 是左单位元且 \(e'\) 是 \(x\) 的右单位元。证明：若 \(\mathbf{G}_{e, e}\) 非空，则它在由 \(\mathbf{E}\) 上的律诱导的律下是一个群。
\end{enumerate}

E若还满足以下条件，则称为群胚：

\begin{enumerate}
\def\labelenumi{(\arabic{enumi})}
\setcounter{enumi}{3}
\tightlist
\item
  对于所有幂等元 \(e, e', G_{e, e'}\) 是非空的。
\end{enumerate}

\begin{enumerate}
\def\labelenumi{(\alph{enumi})}
\setcounter{enumi}{4}
\item
  在群胚 E 中，若 \(a \in G_{e,e'}\) ，证明 \(x \mapsto xa\) 是 \(G_{e,e}\) 到 \(G_{e,e'}, y \mapsto ay\) 上的一个双射映射，是 \(G_{e',e'}\) 到 \(G_{e,e'}\) 上的一个双射映射，而 \(x \mapsto a^{-1}xa\) 是群 \(G_{e,e}\) 到群 \(G_{e',e'}\) 上的一个同构。
\item
  证明 § 1，习题 4(b) 中定义的律在集合 \(\Psi\) 上确定一个群胚结构，如果族 \(\mathfrak{F}\) 的所有集合都是等势的。
\end{enumerate}

\begin{enumerate}
\def\labelenumi{\arabic{enumi}.}
\setcounter{enumi}{23}
\tightlist
\item
  \begin{enumerate}
  \def\labelenumii{(\alph{enumii})}
  \tightlist
  \item
  给定任意集合 E，在集合 \(E \times E\) 上考虑一个并非处处有定义、以乘法记号书写的合成律，在该律下， \((x, y)\) 与 \((y', z)\) 的合成仅当 \(y' = y\) 时才有定义，且此时其值为 \((x, z)\) 。证明 \(\mathbf{E} \times \mathbf{E}\) 配备该合成律是一个群胚（习题 23）。
  \end{enumerate}
\end{enumerate}

\begin{enumerate}
\def\labelenumi{(\alph{enumi})}
\setcounter{enumi}{1}
\tightlist
\item
  设 \((x, y) \equiv (x', y')\) (R) 是一个与上述合成相容的等价关系（即 \((x, y) \equiv (x', y')\) 与 \((y, z) \equiv (y', z')\) 蕴含 \((x, z) \equiv (x', z'))\) ；进一步假设关系 R 满足以下条件：对所有 \(x, y, z\) ，存在且仅存在一个 \(t \in \mathbf{E}\) 使得 \((x, y) \equiv (z, t)\) ，且存在一个 \(u \in \mathbf{E}\) 使得 \((x, y) \equiv (u, z)\) 。
\end{enumerate}

证明在这些条件下， \(E \times E\) 上的群胚结构关于 R 的商结构是一个群结构（先证明商律处处有定义，因为若 \(\dot{x}, \dot{y}, \dot{z}\) 是三个类，满足 \(\dot{x}\dot{y} = \dot{x}\dot{z}\) ，则 \(\dot{y} = \dot{z}\) ；最后建立：对所有 \(x \in E, y \in E, (x, x) \equiv (y, y)\) ）。

\begin{enumerate}
\def\labelenumi{(\alph{enumi})}
\setcounter{enumi}{2}
\tightlist
\item
  设 G 是通过把上述结构赋予 \(E \times E\) 关于 R 的商而得到的群。设 a 为 E 的任意元素；若对所有 \(x \in E\) ， \(f_{a}(x)\) 表示 \((a, x)\) 的模 R 的类，证明 \(f_{a}\) 是 E 到 G 上的一个双射映射，且关系 \((x, y) \equiv (x', y')\) 等价于关系
\end{enumerate}

\[
f _ {a} (x) \left(f _ {a} \left(x ^ {\prime}\right)\right) ^ {- 1} = f _ {a} (y) \left(f _ {a} \left(y ^ {\prime}\right)\right) ^ {- 1}.
\]

要使 G 是交换群，必要且充分的是关系 \((x,y)\equiv (x',y')\) 蕴含 \((x,x')\equiv (y,y')\) 。

¶ 25. 设 E 是一个集合，f 是 \(E^{m}\) 到 E 的映射；我们记

\[
f (x _ {1}, x _ {2}, \dots , x _ {m}) = x _ {1} x _ {2} \dots x _ {m};
\]

假设 \(f\) 满足以下条件：

\[
(x _ {1} x _ {2} \dots x _ {m}) x _ {m + 1} \dots x _ {2 m - 1} = x _ {1} (x _ {2} x _ {3} \dots x _ {m + 1}) x _ {m + 2} \dots x _ {2 m - 1}\tag{1}
\]

恒成立；

\begin{enumerate}
\def\labelenumi{(\arabic{enumi})}
\setcounter{enumi}{1}
\tightlist
\item
  对于所有 \(a_1, a_2, \ldots, a_{m-1}\) ，映射
\end{enumerate}

\[
x \mapsto x a _ {1} a _ {2} \dots a _ {m - 1}
\]

\[
x \mapsto a _ {1} a _ {2} \dots a _ {m - 1} x
\]

是 E 到 E 上的双射。

\begin{enumerate}
\def\labelenumi{(\alph{enumi})}
\tightlist
\item
  证明恒有
\end{enumerate}

\[
(x _ {1} x _ {2} \dots x _ {m}) x _ {m + 1} \dots x _ {2 m - 1} = x _ {1} x _ {2} \dots x _ {i} (x _ {i + 1} \dots x _ {i + m}) x _ {i + m + 1} \dots x _ {2 m - 1}
\]

对每个指标 \(i\) （ \(1 \leqslant i \leqslant m - 1\) ）成立（对 \(i\) 作归纳论证，考虑元素

\[
\left(\left(x _ {1} x _ {2} \dots x _ {m}\right) x _ {m + 1} \dots x _ {2 m - 1}\right) a _ {1} a _ {2} \dots a _ {m - 1}).
\]

\begin{enumerate}
\def\labelenumi{(\alph{enumi})}
\setcounter{enumi}{1}
\tightlist
\item
  对于所有 \(a_1, a_2, \ldots, a_{m-2}\) ，存在 \(u \in E\) 使得对所有 \(x\)
\end{enumerate}

\[
x = x a _ {1} a _ {2} \dots a _ {m - 2} u = u a _ {1} a _ {2} \dots a _ {m - 2} x.
\]

\begin{enumerate}
\def\labelenumi{(\alph{enumi})}
\setcounter{enumi}{2}
\tightlist
\item
  在集合 \(\mathbf{E}^k\) ——它由诸序列 \((u_1, u_2, \ldots, u_k)\) （ \(k\) 个元素，属于 \(\mathbf{E}\) ）组成——（ \(1 \leqslant k \leqslant m - 1\) ）中，考虑表述如下的等价关系 \(\mathbf{R}_k\) ：对所有 \(x_1, x_2, \ldots, x_{m-k}, u_1 u_2 \ldots u_k x_1 x_2 \ldots x_{m-k} = v_1 v_2 \ldots v_k x_1 x_2 \ldots x_{m-k}\) ；令 \(\mathbf{E}_k\) 表示商集 \(\mathbf{E}^k / \mathbf{R}_k\) ，令 \(\mathbf{G}\) 表示 \(\mathbf{E}_1 = \mathbf{E}, \mathbf{E}_2, \ldots, \mathbf{E}_{m-1}\) 的和集（集合论，II，§4，第 8 号）。设 \(\alpha \in \mathbf{E}_i\) , \(\beta \in \mathbf{E}_j\) ；若 \((u_1, u_2, \ldots, u_i)\) 是类 \(\alpha\) 的一个序列，而 \((v_1, \ldots, v_j)\) 是类 \(\beta\) 的一个序列，考虑序列
\end{enumerate}

\[
\left(u _ {1}, u _ {2}, \dots , u _ {i}, v _ {1}, \dots , v _ {j}\right)
\]

（在 \(\mathbf{E}^{i + j}\) 中，若 \(i + j < m\) ），或者序列

\[
\left(u _ {1}, u _ {2}, \dots , u _ {i + j - m}, \left(u _ {i + j - m + 1} \dots u _ {i} v _ {1} v _ {2} \dots v _ {j}\right)\right)
\]

（在 \(E^{i+j-m+1}\) 中，若 \(i+j\geqslant m\) ）；证明该序列在 \(E_{i+j}\) （相应地 \(E_{i+j-m+1}\) ）中的类仅依赖于类 \(\alpha\) 与 \(\beta\) ；把它记作 \(\alpha.\beta\) 。证明这样就在 G 上定义了一个群律， \(H=E_{m-1}\) 是 G 的一个正规子群，且商群 G/H 是 m-1 阶的循环群；最后证明 E 与生成 G/H 的一个类（模 H ）相同，并且 \(x_{1}x_{2}\ldots x_{m}\) 正是序列 \((x_{1},x_{2},\ldots,x_{m})\) 在群 G 中的合成。

¶ 26. 设 G, G' 是两个群， \(f: G \to G'\) 是一个映射，使得对 G 的任意两个元素 \(x, y\) ， \(f(xy) = f(x)f(y)\) 或 \(f(xy) = f(y)f(x)\) 。我们打算证明： \(f\) 是 G 到 G' 中的一个同态，或是 G 到反群 G'0 中的一个同态（换言之，或者 \(f(xy) = f(x)f(y)\) 对每个有序对 \((x, y)\) 成立，或者 \(f(xy) = f(y)f(x)\) 对每个有序对 \((x, y)\) 成立）。

\begin{enumerate}
\def\labelenumi{(\alph{enumi})}
\item
  证明集合 \(N = f(e')\) （ \(e'\) 是 \(G'\) 的单位元）是 G 的一个正规子群，并且 f 通过 \(G \xrightarrow{p} G/N \xrightarrow{g} G'\) 分解，其中 g 是单射。因此可以只考虑 f 为单射的情形，下文中将始终作此假定。
\item
  证明：如果 \(xy = yx\) 那么 \(f(x)f(y) = f(y)f(x)\) （考虑 \(f(x^2y)\) 和 \(f(x^2y^2)\) ，以多种方式表达它们）。
\item
  证明：如果 \(f(xy) = f(x)f(y)\) 那么 \(f(yx) = f(y)f(x)\) 。（借助 (b) 证明，可以只考虑 \(xy \neq yx\) 的情形。）
\item
  证明 \(f(xyx) = \hat{f}(x)f(y)f(x)\) 。（依 \(xy = yx\) 或 \(xy \neq yx\) 而使用 (a) 或 (b)；在第二种情形，考虑 \(f(x^2y)\) 与 \(f(yx^2)\) 。）
\item
  设 A 为这样的 \(x \in G\) 所成的集合：\(f(xy) = f(x)f(y)\) 对所有 \(y \in G\) 成立；设 B 为这样的 \(x \in G\) 所成的集合：\(f(xy) = f(y)f(x)\) 对所有 \(y \in G\) 成立。证明 \(A \neq G\) 、 \(B \neq G\) 与 \(A \cup B = G\) 是不可能的情形。（根据 (b)，此时将存在 A 中的 \(x, y\) ，使得 \(f(xy) = f(x)f(y) \neq f(y)f(x)\) ，以及 B 中的 \(u, v\) ，使得 \(f(uv) = f(v)f(u) \neq f(u)f(v)\) 。依次考虑 \(xu \in A\) 、 \(xu \in B\) 两种可能性，由此推出矛盾；第一种情形考虑 \(f(xuv)\) ，第二种情形考虑 \(f(yxu)\) 。）
\end{enumerate}

证明的其余部分在于用归谬法证明 \(\mathbf{A} \cup \mathbf{B} \neq \mathbf{G}\) 是不可能的。若 \(\mathbf{A} \cup \mathbf{B} \neq \mathbf{G}\) ，则存在 \(a, b, c\) （属于 \(\mathbf{G}\) ），使得

\[
f (a b) = f (a) f (b) \neq f (b) f (a) \quad \text { and } \quad f (a c) = f (c) f (a) \neq f (a) f (c).
\]

\begin{enumerate}
\def\labelenumi{(\alph{enumi})}
\setcounter{enumi}{5}
\item
  证明 \(f(c)f(a)f(b) = f(b)f(a)f(c)\) 。（依 \(bc \neq cb\) 或 \(bc = cb\) 分两种情形考虑。第一种情形考虑 \(f(bac)\) ；第二种情形利用 (b) 与 (c)，并考虑 \(f(abc), f(bca)\) 与 \(f(bac)\) 。）
\item
  考虑 \(f(abac)\) 并得出矛盾。
\end{enumerate}

